\chapter{不确定性传播与置信区域}

\section{协方差近似}
在最优点 $\hat{\bm p}$ 附近，令加权残差向量为 $\bm r(\bm p)$，雅可比为 $\bm J(\bm p)$，则
\[
\bm A=\bm J^{\T}\bm J
\]
可视为信息矩阵近似。若模型线性化误差较小，估计协方差可近似为
\begin{equation}
\bm\Sigma\approx \bm A^{-1}.
\label{eq:cov_approx}
\end{equation}
在牛顿法中也可使用海森矩阵逆近似协方差。两者在残差较小时通常等价到同一量级。

\section{置信椭球构造}
定义二次型
\[
Q(\bm p)=
(\bm p-\hat{\bm p})^{\T}\bm\Sigma^{-1}(\bm p-\hat{\bm p}).
\]
给定置信水平 $p$，令 $\chi^2_{3,p}$ 为自由度 3 的卡方分位数，则置信区域可写为
\begin{equation}
\mathcal E_p=
\left\{\bm p\in\R^3:\ Q(\bm p)\le \chi^2_{3,p}\right\}.
\end{equation}
将 $\bm\Sigma$ 作特征分解
\[
\bm\Sigma=\bm U\diag(\lambda_1,\lambda_2,\lambda_3)\bm U^{\T},
\]
则椭球三条半轴长度为
\[
a_i=\sqrt{\chi^2_{3,p}\lambda_i},\quad i=1,2,3.
\]

\section{与代码输出的一致性}
代码中基于目标函数值与海森矩阵特征值给出 95\% 和 99\% 置信椭球半轴。该输出与上节理论形式一致，区别仅在于符号缩放方式和自由度取值设置。工程应用中建议明确自由度定义并统一统计口径，以便跨项目比较。

\section{主轴方向解释}
设特征向量矩阵 $\bm U=[\bm u_1,\bm u_2,\bm u_3]$。则：
\begin{itemize}
  \item $\bm u_i$ 给出不确定性的主方向；
  \item $\lambda_i$ 越大，沿 $\bm u_i$ 方向方差越大；
  \item 若第三主轴对应近竖直方向且半轴显著更长，说明高程估计较弱。
\end{itemize}
这与第 6 章误差统计结论保持一致。

\section{误差传播分析}
考虑观测噪声微扰 $\delta\bm y$，由线性化关系有
\[
\delta\bm p\approx(\bm J^{\T}\bm W\bm J)^{-1}\bm J^{\T}\bm W\,\delta\bm y.
\]
因此误差传播受三方面影响：
\begin{enumerate}
  \item 观测噪声本身（$\delta\bm y$）；
  \item 权重设计（$\bm W$）；
  \item 几何条件（$\bm J$ 的列相关性）。
\end{enumerate}
几何退化时，即使噪声很小也可能导致较大位置误差。

\section{GDOP 类指标}
为了量化几何优劣，可定义
\[
GDOP=\sqrt{\operatorname{trace}(\bm\Sigma)}.
\]
并进一步分解为
\[
HDOP=\sqrt{\Sigma_{xx}+\Sigma_{yy}},\qquad
VDOP=\sqrt{\Sigma_{zz}}.
\]
在本文实验中，$VDOP$ 通常显著大于 $HDOP$，反映垂向几何弱化现象。

\section{稳健置信评估建议}
在实际工程中，噪声可能偏离高斯假设。建议：
\begin{itemize}
  \item 采用残差归一化检验，筛查异常观测；
  \item 使用重采样（Bootstrap）估计经验置信区间；
  \item 对非高斯噪声引入 M-估计或 Huber 损失函数。
\end{itemize}

\section{本章小结}
本章建立了从矩阵近似到置信椭球构造的完整不确定性评估流程，并解释了估计误差方向性来源。该部分使定位结果从“点估计”提升为“带可信度的估计”。
